\section{本章小结}
\label{sec:zo-summary}

本章在只能读函数值的条件下回答了怎样构造方向、怎样迭代、最后得到什么保证。函数值反馈先要说清接口：精确值、加性噪声和样本函数是三种读数模型，单点与成对、独立噪声与共享样本是不同的读取条件，算法抽方向的随机性与读数的随机性分开记录。

局部方向有两种读法。有限差分以原梯度为对象：坐标差分需要$n+1$或$2n$次读数，随机方向差分和SPSA只需两次，但一条方向只有一条方向的信息，有限扰动下一般有偏，SPSA的坐标尺度乘以$\sqrt{n}$才是位移长度。随机平滑以平滑目标为对象：球平滑的恒等式~\eqref{eq:zo-ball-identity}和高斯平滑的恒等式~\eqref{eq:zo-gauss-identity}给出严格无偏的估计，单点估计受函数值大小影响，双点估计消去常数偏移。独立噪声让差商方差按$1/\delta^2$增长，共享样本是否改善取决于差值的波动。

把估计器代入迭代后，算法实际优化的目标由条件均值决定。步长受维数限制约为$1/(nL)$，扰动尺度决定平滑目标和噪声，缩小它要同时记入目标变化和方差增长。约束下要分清更新可行与查询可行，球形扰动可用收缩集合，高斯扰动有无界支撑。输出按问题选择平均点或随机返回点，预算耗尽、读数不变和驻点保证是不同的停止依据，含噪比较需要独立复评。

按读数计费时，两次查询并不代表成本与维数无关，因子$n$转移到了步长上。平滑目标上的结论要加上平滑误差才是原目标的结论，原目标不光滑时平滑梯度小只是较弱的驻点概念，也不保证全局最优。选择方法时，先看接口：能读坐标且读数便宜就用差分；扰动必须同时施加时用SPSA；只有单点反馈时用单点估计；能共享样本时优先双点估计并尽量取小$\delta$。

高斯平滑的恒等式也可以看作：把$\bm{x}$当作高斯搜索分布的均值，这个估计描述均值移动时平均目标的变化。第~\ref{chap:random-search}章继续讨论怎样更新搜索分布和种群，那里的方法不都等同于本章的零阶梯度下降。
